% =============================================================================
\section{Polarización y densidades de carga ligada}\label{sec:polarizacion}
% =============================================================================

\subsection{El vector polarización}\label{sec:vector-P}

\begin{definicion}[Vector polarización]\label{def:polarizacion}
Se llama \emph{vector polarización} del dieléctrico en el punto $\vr\,'$ al momento dipolar por unidad de
volumen,
\begin{equation}\label{eq:P}
  \vP(\vr\,')\defeq\frac{\dd\vp}{\dd V'} ,
\end{equation}
donde $\dd\vp$ es el momento dipolar neto del elemento $\dd V'$. Sus unidades son $\si{C\,m/m^3}=\si{C/m^2}$. Si hay
$N$ moléculas por unidad de volumen con momento medio $\langle\vp\rangle$, $\vP=N\langle\vp\rangle$. Se dice que el
dieléctrico está \emph{uniformemente polarizado} si $\vP$ no depende del punto.
\end{definicion}

Con esta definición, el momento dipolar de cada elemento es $\dd\vp=\vP\,\dd V'$, y el potencial del
dieléctrico completo se obtiene sumando \eqref{eq:dphi-dipolo}:
\begin{equation}\label{eq:phi-P}
  \phi(\vr)=\frac{1}{4\pi\eo}\int_{V'}\frac{\vP(\vr\,')\cdot\uv R}{R^2}\,\dd V' .
\end{equation}
Esta expresión es correcta, pero poco transparente. El resultado central del tema es que puede
reescribirse como el potencial de una distribución de carga ordinaria.

\subsection{Densidades de carga de polarización}\label{sec:densidades}

\begin{lema}[Gradiente de $1/R$]\label{lem:grad-1R}
Para $\vr\neq\vr\,'$,
\begin{equation}\label{eq:grad-1R}
  \grad\Big(\frac1R\Big)=-\frac{\uv R}{R^2},\qquad
  \grad'\Big(\frac1R\Big)=+\frac{\uv R}{R^2}.
\end{equation}
\end{lema}

\begin{proof}
Para cualquier $n$, $\grad R^n=nR^{n-1}\grad R=nR^{n-1}\uv R$, pues $\grad R=\vR/R$~\cite{T0}. Con $n=-1$ se
obtiene la primera igualdad. Como $R$ depende de $\vr-\vr\,'$, derivar respecto a $\vr\,'$ cambia el signo,
$\grad'f(\vr-\vr\,')=-\grad f(\vr-\vr\,')$, de donde la segunda.
\end{proof}

\begin{teorema}[Densidades de carga de polarización]\label{teo:densidades}
El potencial \eqref{eq:phi-P} de un dieléctrico polarizado que ocupa la región $V'$, de frontera
$S'=\partial V'$ con normal exterior $\nn$, es el de una densidad volúmica $\rhoP$ y una densidad superficial
$\sigmaP$ de carga,
\begin{equation}\label{eq:phi-densidades}
  \phi(\vr)=\frac{1}{4\pi\eo}\int_{V'}\frac{\rhoP(\vr\,')}{R}\,\dd V'
            +\frac{1}{4\pi\eo}\oiint_{S'}\frac{\sigmaP(\vr\,')}{R}\,\dd S' ,
\end{equation}
con
\begin{equation}\label{eq:rhoP-sigmaP}
  \boxed{\;\rhoP(\vr\,')=-\grad'\!\cdot\vP(\vr\,'),\qquad \sigmaP(\vr\,')=\vP(\vr\,')\cdot\nn .\;}
\end{equation}
Se llaman \emph{densidades de carga de polarización} (o de carga ligada).
\end{teorema}

\begin{proof}
Por el \cref{lem:grad-1R}, el integrando de \eqref{eq:phi-P} es $\vP\cdot\grad'(1/R)$. La identidad del producto
$\grad'\!\cdot(f\vA)=f\,\grad'\!\cdot\vA+\vA\cdot\grad'f$, con $f=1/R$ y $\vA=\vP$, da
\[
  \frac{\vP\cdot\uv R}{R^2}=\vP\cdot\grad'\Big(\frac1R\Big)
  =-\frac{\grad'\!\cdot\vP}{R}+\grad'\!\cdot\Big(\frac{\vP}{R}\Big).
\]
Integrando en $V'$ y usando la linealidad de la integral,
\[
  \phi(\vr)=\frac{1}{4\pi\eo}\left[\int_{V'}\frac{-\grad'\!\cdot\vP}{R}\,\dd V'
            +\int_{V'}\grad'\!\cdot\Big(\frac{\vP}{R}\Big)\dd V'\right].
\]
A la segunda integral se le aplica el teorema de la divergencia de Gauss,
$\iiint_{V}\grad\cdot\vA\,\dd V=\oiint_{\partial V}\vA\cdot\dS$, con $\dS'=\nn\,\dd S'$:
\[
  \iiint_{V'}\grad'\!\cdot\Big(\frac{\vP}{R}\Big)\dd V'=\oiint_{S'}\frac{\vP\cdot\nn}{R}\,\dd S' .
\]
Basta llamar $\rhoP$ y $\sigmaP$ a los numeradores para obtener \eqref{eq:phi-densidades}.
\end{proof}

\begin{observacion}[Sobre el rigor de la deducción]\label{obs:singular}
Si $\vr$ está dentro del dieléctrico, el integrando de \eqref{eq:phi-P} es singular en $\vr\,'=\vr$, pero solo como
$1/R^2$, que es integrable en tres dimensiones; basta excluir una bola de radio $\delta$ alrededor de $\vr$ y hacer
$\delta\to0$, pues la contribución de su superficie es $\mathcal O(\delta)$. Por otra parte, la aproximación dipolar
promedia sobre volúmenes que contienen muchas moléculas: los campos que obtenemos son \emph{macroscópicos}, y en
particular \eqref{eq:phi-densidades} no describe las fluctuaciones a escala atómica.
\end{observacion}

\begin{corolario}[Las cargas ligadas son cargas de verdad]\label{cor:campo-ligadas}
El campo de un dieléctrico polarizado es $\vE=-\grad\phi$, es decir,
\begin{equation}\label{eq:E-densidades}
  \vE(\vr)=\frac{1}{4\pi\eo}\int_{V'}\frac{\rhoP\,\uv R}{R^2}\,\dd V'
           +\frac{1}{4\pi\eo}\oiint_{S'}\frac{\sigmaP\,\uv R}{R^2}\,\dd S' ,
\end{equation}
el mismo que crearían cargas libres con esas densidades. Toda la electrostática del vacío se aplica a ellas.
\end{corolario}

\begin{proof}
Se deriva \eqref{eq:phi-densidades} bajo el signo integral usando la primera igualdad de \eqref{eq:grad-1R}.
\end{proof}

\subsection{Interpretación física}\label{sec:interpretacion}

¿De dónde salen estas cargas, si el dieléctrico es neutro? Pensemos en una cadena de dipolos idénticos
(\cref{fig:cadena}\,(a)): la carga positiva de cada dipolo compensa la negativa del siguiente, y solo
«sobran» la carga negativa del primero y la positiva del último. La cadena equivale a dos cargas $\mp q$ en
sus extremos.

% Figura: interpretación de las densidades de polarización. Cadena de dipolos,
% polarización uniforme (solo carga superficial) y no uniforme (carga de volumen).
\begin{figure}[t]
\centering
\begin{subfigure}[b]{\textwidth}\centering
\begin{tikzpicture}[scale=1.0]
  \foreach \i in {0,...,4} {
    \pgfmathsetmacro{\x}{1.05*\i}
    \draw[campoP, line width=0.9pt, -{Stealth[length=1.6mm]}] (\x,0) -- (\x+0.7,0);
    \node[cneg] at (\x+0.02,-0.27) {$-$};
    \node[cpos] at (\x+0.68,-0.27) {$+$};
  }
  \draw[gauss, rounded corners=3pt] (1.65,-0.48) rectangle (2.20,0.3);
  \node[etiqt, anchor=north] at (1.93,-0.55) {neutro};
  \node[etiq] at (5.75,-0.05) {$=$};
  \draw[texto, line width=0.6pt] (6.4,0) -- (10.5,0);
  \fill[carganeg] (6.4,0) circle (0.08); \fill[cargapos] (10.5,0) circle (0.08);
  \node[etiq, anchor=north] at (6.4,-0.12) {$-q$};
  \node[etiq, anchor=north] at (10.5,-0.12) {$+q$};
\end{tikzpicture}
\caption{Cadena de dipolos idénticos: el interior es localmente neutro y solo «sobran» las cargas de los extremos}
\end{subfigure}\\[0.9em]
\begin{subfigure}[b]{0.48\textwidth}\centering
\begin{tikzpicture}[scale=0.95]
  \fill[dielec] (0,0) rectangle (3.4,2.0);
  \draw[borde] (0,0) rectangle (3.4,2.0);
  \foreach \i in {0,...,4} \foreach \j in {0,...,4}
    \draw[campoP, line width=0.75pt, -{Stealth[length=1.5mm,width=1.2mm]}]
      ({0.2+0.65*\i},{0.22+0.39*\j}) -- ({0.62+0.65*\i},{0.22+0.39*\j});
  \foreach \j in {0,...,4} {
    \node[cneg] at (-0.18,{0.22+0.39*\j}) {$-$};
    \node[cpos] at (3.58,{0.22+0.39*\j}) {$+$};
  }
  \draw[cargapos, dashed, line width=0.7pt] (1.72,1.0) ellipse (0.35 and 0.8);
  \draw[vecn] (0,1.0) -- (-0.75,1.0) node[above, etiq] {$\nn$};
  \draw[vecn] (3.4,1.0) -- (4.15,1.0) node[above, etiq] {$\nn$};
\end{tikzpicture}
\caption{$\vP$ uniforme: $\rhoP=0$, $\sigmaP=\vP\cdot\nn$}
\end{subfigure}\hfill
\begin{subfigure}[b]{0.48\textwidth}\centering
\begin{tikzpicture}[scale=0.95]
  \fill[dielec] (0,0) rectangle (3.4,2.0);
  \draw[borde] (0,0) rectangle (3.4,2.0);
  \foreach \i/\L in {0/0.18,1/0.28,2/0.38,3/0.48,4/0.58} \foreach \j in {0,...,4}
    \draw[campoP, line width=0.75pt, -{Stealth[length=1.5mm,width=1.2mm]}]
      ({0.42+0.65*\i-\L/2},{0.22+0.39*\j}) -- ({0.42+0.65*\i+\L/2},{0.22+0.39*\j});
  \foreach \j in {0,...,4} {
    \node[cneg] at (-0.18,{0.22+0.39*\j}) {$-$};
    \node[cpos] at (3.58,{0.22+0.39*\j}) {$+$};
    \node[cpos] at (3.78,{0.22+0.39*\j}) {$+$};
  }
  \draw[cargapos, dashed, line width=0.7pt] (1.72,1.0) ellipse (0.35 and 0.8);
  \draw[vecn] (0,1.0) -- (-0.75,1.0) node[above, etiq] {$\nn$};
  \draw[vecn] (3.9,1.0) -- (4.55,1.0) node[above, etiq] {$\nn$};
\end{tikzpicture}
\caption{$\vP$ no uniforme: $\rhoP=-\divg\vP\neq0$}
\end{subfigure}
\caption{Origen de las cargas de polarización. (a) En una cadena de dipolos la carga positiva de cada
uno compensa la negativa del siguiente. (b) Si $\vP$ es uniforme, un volumen interior cualquiera (trazos)
contiene tantas cabezas como colas de dipolos y es neutro: la carga ligada aparece solo en la superficie.
(c) Si $\vP$ crece hacia la derecha, en el mismo volumen entran más extremos negativos que positivos
salen y queda una carga neta: es la densidad volúmica $\rhoP=-\divg\vP$.}
\label{fig:cadena}
\end{figure}


Lo mismo ocurre en un bloque uniformemente polarizado (\cref{fig:cadena}\,(b)): cualquier volumen interior
contiene tantas cabezas como colas de dipolos y es neutro, y solo aparece carga en las caras por las que
$\vP$ entra (carga negativa) o sale (positiva). Si la polarización no es uniforme
(\cref{fig:cadena}\,(c)), en un volumen interior entran y salen distintos números de extremos, y queda una
carga neta: esa es la densidad $\rhoP$.

\begin{ejemplo}[Cálculo directo de $\sigmaP$]\label{ej:sigmaP-directo}
Sea un cilindro de longitud $\ell$ y sección $A$ uniformemente polarizado según su eje
(\cref{fig:cilindro-uniforme}\,(a)). Su momento dipolar es $p=P\,A\ell$, y el de un par de cargas $\mp q$ en las
caras es $p=q\ell$. Igualando, $q=PA$, de modo que la densidad en cada cara es $\sigmaP=q/A=P$. Si la cara
se corta oblicuamente, con normal $\nn$ formando un ángulo $\theta$ con $\vP$ (\cref{fig:cilindro-uniforme}\,(b)), la
misma carga se reparte en un área $A'=A/\cos\theta$, y
\[
  \sigmaP=\frac{q}{A'}=\frac{PA}{A/\cos\theta}=P\cos\theta=\vP\cdot\nn ,
\]
que es exactamente la expresión del \cref{teo:densidades}. En la superficie lateral $\vP\perp\nn$ y $\sigmaP=0$.
\end{ejemplo}

% Figura: cálculo directo de sigma_P en un cilindro uniformemente polarizado,
% corte oblicuo (sigma_P = P cos theta) y polarización radial (rho_P < 0).
\begin{figure}[t]
\centering
\begin{subfigure}[b]{0.46\textwidth}\centering
\begin{tikzpicture}[scale=0.9]
  % cilindro (vista lateral con perspectiva caballera)
  \fill[dielec] (0,-0.6) -- (2.6,-0.6) arc (-90:90:0.24 and 0.6) -- (0,0.6) arc (90:270:0.24 and 0.6);
  \draw[borde] (0,-0.6) -- (2.6,-0.6) (0,0.6) -- (2.6,0.6);
  \draw[borde] (0,0) ellipse (0.24 and 0.6);
  \fill[dielecB] (2.6,0) ellipse (0.24 and 0.6);
  \draw[borde] (2.6,0) ellipse (0.24 and 0.6);
  \node[etiq] at (2.6,0) {$A$};
  \draw[vecP] (0.6,0) -- (1.9,0) node[above, midway, etiq] {$\vP$};
  \draw[cota] (0,0.92) -- (2.6,0.92) node[midway, above, etiq] {$\ell$};
  % equivalente
  \node[etiq] at (3.45,0) {$\equiv$};
  \draw[borde] (4.3,0) ellipse (0.24 and 0.6);
  \draw[borde] (6.6,0) ellipse (0.24 and 0.6);
  \draw[texto!70!fondo, line width=0.5pt] (4.3,0) -- (6.6,0);
  \node[cneg, font=\small] at (4.3,0.85) {$-q$};
  \node[cpos, font=\small] at (6.6,0.85) {$+q$};
  \draw[cota] (4.3,-0.92) -- (6.6,-0.92) node[midway, below, etiq] {$\ell$};
\end{tikzpicture}
\caption{$p=q\ell=P\,A\ell\ \Rightarrow\ \sigmaP=q/A=P$}
\end{subfigure}\hfill
\begin{subfigure}[b]{0.27\textwidth}\centering
\begin{tikzpicture}[scale=0.9]
  \fill[dielec] (0,-0.6) -- (2.2,-0.6) -- (2.75,0.6) -- (0,0.6) arc (90:270:0.24 and 0.6);
  \draw[borde] (0,-0.6) -- (2.2,-0.6) (0,0.6) -- (2.75,0.6);
  \draw[borde] (0,0) ellipse (0.24 and 0.6);
  \fill[dielecB] (2.475,0) ellipse [x radius=0.24, y radius=0.66, rotate=-24.6];
  \draw[borde] (2.475,0) ellipse [x radius=0.24, y radius=0.66, rotate=-24.6];
  \node[etiq] at (3.05,0.82) {$A'$};
  \draw[vecP] (0.45,0) -- (1.65,0) node[above, midway, etiq] {$\vP$};
  \draw[vecn] (2.475,0) -- ++(-24.6:1.05) node[right, etiq] {$\nn$};
  \draw[texto!60!fondo, dashed, line width=0.35pt] (2.475,0) -- (3.75,0);
  \draw[tenue, line width=0.4pt] (3.375,0) arc (0:-24.6:0.9);
  \node[etiqt] at ($(2.475,0)+(-11:1.1)$) {$\theta$};
\end{tikzpicture}
\caption{$\sigmaP=\dfrac{q}{A'}=P\cos\theta$}
\end{subfigure}\hfill
\begin{subfigure}[b]{0.23\textwidth}\centering
\begin{tikzpicture}[scale=0.9]
  \fill[dielec] (0,0) circle (1.05);
  \draw[equipot] (0,0) circle (1.05);
  \foreach \t in {0,45,...,315} {
    \draw[campoP, line width=0.8pt, -{Stealth[length=1.6mm]}] (\t:0.3) -- (\t:0.95);
    \node[cpos] at (\t:1.28) {$+$};
  }
  \foreach \t in {22.5,112.5,202.5,292.5} \node[cneg, font=\small\bfseries] at (\t:0.15) {$-$};
\end{tikzpicture}
\caption{$\divg\vP>0$}
\end{subfigure}
\caption{(a) Un cilindro uniformemente polarizado equivale a dos discos con cargas $\mp q$: de
$P=p/V=q\ell/(A\ell)$ se sigue $\sigmaP=P$. (b) Si la cara se corta oblicuamente, la misma carga se
reparte sobre un área mayor $A'=A/\cos\theta$, y $\sigmaP=P\cos\theta=\vP\cdot\nn$. (c) Una polarización
divergente arrastra carga positiva hacia la superficie y deja un defecto negativo en el interior:
$\rhoP=-\divg\vP<0$ y $\sigmaP>0$, con carga ligada total nula.}
\label{fig:cilindro-uniforme}
\end{figure}


\subsection{La carga de polarización total es nula}\label{sec:carga-total}

Una polarización divergente arrastra carga positiva hacia la superficie y deja un defecto negativo en el
interior (\cref{fig:cilindro-uniforme}\,(c)): la carga que falta en el volumen aparece en la superficie,
$\int_{V'}\rhoP\,\dd V'=-\oiint_{S'}\sigmaP\,\dd S'$. Es una consecuencia general.

\begin{proposicion}[Neutralidad de la carga ligada]\label{prop:QP-nula}
La carga de polarización total de un dieléctrico es nula:
\begin{equation}\label{eq:QP}
  \QP=\int_{V'}\rhoP\,\dd V'+\oiint_{S'}\sigmaP\,\dd S'=0 .
\end{equation}
\end{proposicion}

\begin{proof}
Por el teorema de la divergencia,
$\int_{V'}\rhoP\,\dd V'=-\int_{V'}\grad'\!\cdot\vP\,\dd V'=-\oiint_{S'}\vP\cdot\nn\,\dd S'=-\oiint_{S'}\sigmaP\,\dd S'$.
\end{proof}

\begin{observacion}[El momento dipolar de las cargas ligadas]\label{obs:momento-ligado}
Como el monopolo de las cargas ligadas es nulo, el término dominante de su desarrollo multipolar~\cite{T2} es
el dipolar, y su momento es el esperado:
\[
  \int_{V'}\vr\,'\rhoP\,\dd V'+\oiint_{S'}\vr\,'\sigmaP\,\dd S'=\int_{V'}\vP\,\dd V' .
\]
Para la componente $x$, por ejemplo, $\grad'\!\cdot(x'\vP)=x'\grad'\!\cdot\vP+P_x$, y el teorema de la divergencia da
$\oiint_{S'}x'\vP\cdot\nn\,\dd S'-\int_{V'}x'\grad'\!\cdot\vP\,\dd V'=\int_{V'}P_x\,\dd V'$. Lejos del dieléctrico, su campo es el
de un dipolo $\vp_{\text{tot}}=\int\vP\,\dd V'$, como debe ser.
\end{observacion}

La comprobación explícita de \eqref{eq:QP} en un cubo con polarización radial es el \cref{prob:cubo} de
los seminarios, y en una esfera uniformemente polarizada, el \cref{prob:esfera-polarizada}.
